\chapter{Unit 1: Python Environment Setup, Language Basics, and Computational Foundations}

\section{Section A: Multiple Choice Practice Questions (60 Curated Problems)}

\noindent\textbf{Q1.} Which tool is commonly used to write and run Python code? \hfill [Tag: Installing Python and IDEs (Anaconda, Jupyter, VS Code) | Easy]
\begin{enumerate}[label=(\alph*),topsep=2pt,itemsep=1pt,partopsep=0pt,parsep=0pt]
    \item {JPEG}
    \item {VS Code}
    \item {MP3}
    \item {HTML}
\end{enumerate}
\vspace{2mm}
\noindent\textbf{Q2.} Which symbol starts a comment in Python? \hfill [Tag: Python syntax | Easy]
\begin{enumerate}[label=(\alph*),topsep=2pt,itemsep=1pt,partopsep=0pt,parsep=0pt]
    \item {\#}
    \item {--}
    \item {//}
    \item {/*}
\end{enumerate}
\vspace{2mm}
\noindent\textbf{Q3.} Which statement assigns the value 10 to a variable named score ? \hfill [Tag: Variables | Easy]
\begin{enumerate}[label=(\alph*),topsep=2pt,itemsep=1pt,partopsep=0pt,parsep=0pt]
    \item {score = 10}
    \item {10 = score}
    \item {set score 10}
    \item {score == 10}
\end{enumerate}
\vspace{2mm}
\noindent\textbf{Q4.} What is the data type of 3.14 in Python? \hfill [Tag: Data types | Easy]
\begin{enumerate}[label=(\alph*),topsep=2pt,itemsep=1pt,partopsep=0pt,parsep=0pt]
    \item {float}
    \item {str}
    \item {int}
    \item {bool}
\end{enumerate}
\vspace{2mm}
\noindent\textbf{Q5.} Which operator is used for multiplication in Python? \hfill [Tag: Operators | Easy]
\begin{enumerate}[label=(\alph*),topsep=2pt,itemsep=1pt,partopsep=0pt,parsep=0pt]
    \item {x}
    \item {*}
    \item {\textasciicircum{}}
    \item {\%}
\end{enumerate}
\vspace{2mm}
\noindent\textbf{Q6.} What is the value of the expression 2 + 3 * 4 ? \hfill [Tag: Expressions | Easy]
\begin{enumerate}[label=(\alph*),topsep=2pt,itemsep=1pt,partopsep=0pt,parsep=0pt]
    \item {20}
    \item {9}
    \item {24}
    \item {14}
\end{enumerate}
\vspace{2mm}
\noindent\textbf{Q7.} Which Python function displays text on the screen? \hfill [Tag: Input/output | Easy]
\begin{enumerate}[label=(\alph*),topsep=2pt,itemsep=1pt,partopsep=0pt,parsep=0pt]
    \item {input()}
    \item {display()}
    \item {showtext()}
    \item {print()}
\end{enumerate}
\vspace{2mm}
\noindent\textbf{Q8.} Which keyword is commonly used to bring a module into a Python program? \hfill [Tag: Module basics | Easy]
\begin{enumerate}[label=(\alph*),topsep=2pt,itemsep=1pt,partopsep=0pt,parsep=0pt]
    \item {import}
    \item {using}
    \item {include}
    \item {module}
\end{enumerate}
\vspace{2mm}
\noindent\textbf{Q9.} Which statement imports Python's math module? \hfill [Tag: Math module | Easy]
\begin{enumerate}[label=(\alph*),topsep=2pt,itemsep=1pt,partopsep=0pt,parsep=0pt]
    \item {use math}
    \item {math.load()}
    \item {import math}
    \item {include math}
\end{enumerate}
\vspace{2mm}
\noindent\textbf{Q10.} Which data type is used to represent text in Python? \hfill [Tag: Representing text | Easy]
\begin{enumerate}[label=(\alph*),topsep=2pt,itemsep=1pt,partopsep=0pt,parsep=0pt]
    \item {list}
    \item {int}
    \item {str}
    \item {float}
\end{enumerate}
\vspace{2mm}
\noindent\textbf{Q11.} Which is a valid Python string? \hfill [Tag: String basics | Easy]
\begin{enumerate}[label=(\alph*),topsep=2pt,itemsep=1pt,partopsep=0pt,parsep=0pt]
    \item {"Hello"}
    \item {\{Hello\}}
    \item {Hello}
    \item {[Hello]}
\end{enumerate}
\vspace{2mm}
\noindent\textbf{Q12.} Which symbol is used to create a list in Python? \hfill [Tag: List | Easy]
\begin{enumerate}[label=(\alph*),topsep=2pt,itemsep=1pt,partopsep=0pt,parsep=0pt]
    \item {Curly braces: \{\}}
    \item {Parentheses: ()}
    \item {Angle brackets: $<$$>$}
    \item {Square brackets: []}
\end{enumerate}
\vspace{2mm}
\noindent\textbf{Q13.} Which data type stores either True or False ? \hfill [Tag: Common data types summary | Easy]
\begin{enumerate}[label=(\alph*),topsep=2pt,itemsep=1pt,partopsep=0pt,parsep=0pt]
    \item {str}
    \item {list}
    \item {float}
    \item {bool}
\end{enumerate}
\vspace{2mm}
\noindent\textbf{Q14.} What does int("7") return? \hfill [Tag: Type conversions | Easy]
\begin{enumerate}[label=(\alph*),topsep=2pt,itemsep=1pt,partopsep=0pt,parsep=0pt]
    \item {The string "7"}
    \item {The float 7.0}
    \item {The integer 7}
    \item {The Boolean True}
\end{enumerate}
\vspace{2mm}
\noindent\textbf{Q15.} What decimal number is represented by the binary number 101 ? \hfill [Tag: Binary numbers | Easy]
\begin{enumerate}[label=(\alph*),topsep=2pt,itemsep=1pt,partopsep=0pt,parsep=0pt]
    \item {3}
    \item {6}
    \item {5}
    \item {4}
\end{enumerate}
\vspace{2mm}
\noindent\textbf{Q16.} Which is an f-string in Python? \hfill [Tag: String formatting | Easy]
\begin{enumerate}[label=(\alph*),topsep=2pt,itemsep=1pt,partopsep=0pt,parsep=0pt]
    \item {string("Hello", name)}
    \item {"f Hello name"}
    \item {format Hello name}
    \item {f"Hello \{name\}"}
\end{enumerate}
\vspace{2mm}
\noindent\textbf{Q17.} What is the result of entering 6 / 2 in the Python shell? \hfill [Tag: Using Python shell as a calculator and running simple scripts | Easy]
\begin{enumerate}[label=(\alph*),topsep=2pt,itemsep=1pt,partopsep=0pt,parsep=0pt]
    \item {3.0}
    \item {12}
    \item {4}
    \item {2}
\end{enumerate}
\vspace{2mm}
\noindent\textbf{Q18.} What is PEP 8 mainly concerned with? \hfill [Tag: PEP 8 | Easy]
\begin{enumerate}[label=(\alph*),topsep=2pt,itemsep=1pt,partopsep=0pt,parsep=0pt]
    \item {Managing internet accounts}
    \item {Installing computer hardware}
    \item {Python style guidelines}
    \item {Creating image files}
\end{enumerate}
\vspace{2mm}
\noindent\textbf{Q19.} Which structure chooses between two blocks of code? \hfill [Tag: If-else statement | Easy]
\begin{enumerate}[label=(\alph*),topsep=2pt,itemsep=1pt,partopsep=0pt,parsep=0pt]
    \item {for-in}
    \item {if-else}
    \item {def-return}
    \item {import-from}
\end{enumerate}
\vspace{2mm}
\noindent\textbf{Q20.} Which operator checks whether two values are equal? \hfill [Tag: Equality and relational operators | Easy]
\begin{enumerate}[label=(\alph*),topsep=2pt,itemsep=1pt,partopsep=0pt,parsep=0pt]
    \item {==}
    \item {=$>$}
    \item {!=}
    \item {=}
\end{enumerate}
\vspace{2mm}
\noindent\textbf{Q21.} A student wants an environment that includes Python, commonly used scientific libraries, and a package manager with minimal separate setup. Which choice is most suitable? \hfill [Tag: Installing Python and IDEs (Anaconda, Jupyter, VS Code) | Medium]
\begin{enumerate}[label=(\alph*),topsep=2pt,itemsep=1pt,partopsep=0pt,parsep=0pt]
    \item {A standard Python interpreter only}
    \item {Anaconda distribution}
    \item {A text editor without extensions}
    \item {The Python documentation website}
\end{enumerate}
\vspace{2mm}
\noindent\textbf{Q22.} Which code correctly defines a function that returns the square of its argument? \hfill [Tag: Python syntax | Medium]
\begin{enumerate}[label=(\alph*),topsep=2pt,itemsep=1pt,partopsep=0pt,parsep=0pt]
    \item {define square(x): x ** 2}
    \item {function square(x): return x ** 2}
    \item {def square(x) return x ** 2}
    \item {def square(x): return x ** 2}
\end{enumerate}
\vspace{2mm}
\noindent\textbf{Q23.} After executing a = 5 and a = a + 3 , what value does a contain? \hfill [Tag: Variables | Medium]
\begin{enumerate}[label=(\alph*),topsep=2pt,itemsep=1pt,partopsep=0pt,parsep=0pt]
    \item {8}
    \item {5}
    \item {15}
    \item {2}
\end{enumerate}
\vspace{2mm}
\noindent\textbf{Q24.} What is the data type of the value produced by 7 / 2 in Python 3? \hfill [Tag: Data types | Medium]
\begin{enumerate}[label=(\alph*),topsep=2pt,itemsep=1pt,partopsep=0pt,parsep=0pt]
    \item {str}
    \item {float}
    \item {complex}
    \item {int}
\end{enumerate}
\vspace{2mm}
\noindent\textbf{Q25.} What is the value of 17 // 5 + 17 \% 5 ? \hfill [Tag: Operators | Medium]
\begin{enumerate}[label=(\alph*),topsep=2pt,itemsep=1pt,partopsep=0pt,parsep=0pt]
    \item {7}
    \item {3}
    \item {8}
    \item {5}
\end{enumerate}
\vspace{2mm}
\noindent\textbf{Q26.} What is the result of the expression 2 * (3 + 4) ** 2 ? \hfill [Tag: Expressions | Medium]
\begin{enumerate}[label=(\alph*),topsep=2pt,itemsep=1pt,partopsep=0pt,parsep=0pt]
    \item {fifty-six}
    \item {one hundred}
    \item {one hundred forty-four}
    \item {ninety-eight}
\end{enumerate}
\vspace{2mm}
\noindent\textbf{Q27.} What is displayed by this code if the user enters 12 ? age = input("Age: ") print(age + " years") \hfill [Tag: Input/output | Medium]
\begin{enumerate}[label=(\alph*),topsep=2pt,itemsep=1pt,partopsep=0pt,parsep=0pt]
    \item {A type error occurs}
    \item {12 years}
    \item {age years}
    \item {24 years}
\end{enumerate}
\vspace{2mm}
\noindent\textbf{Q28.} Which statement correctly imports only the sqrt function from the math module? \hfill [Tag: Module basics | Medium]
\begin{enumerate}[label=(\alph*),topsep=2pt,itemsep=1pt,partopsep=0pt,parsep=0pt]
    \item {from math import sqrt}
    \item {import sqrt from math}
    \item {using math.sqrt}
    \item {include math.sqrt}
\end{enumerate}
\vspace{2mm}
\noindent\textbf{Q29.} What value is returned by math.ceil(4.2) after import math ? \hfill [Tag: Math module | Medium]
\begin{enumerate}[label=(\alph*),topsep=2pt,itemsep=1pt,partopsep=0pt,parsep=0pt]
    \item {6}
    \item {4}
    \item {5}
    \item {4.2}
\end{enumerate}
\vspace{2mm}
\noindent\textbf{Q30.} Why does Python support Unicode strings? \hfill [Tag: Representing text | Medium]
\begin{enumerate}[label=(\alph*),topsep=2pt,itemsep=1pt,partopsep=0pt,parsep=0pt]
    \item {To convert numbers automatically into text}
    \item {To prevent strings from being indexed}
    \item {To represent characters from many writing systems}
    \item {To make every string use one byte}
\end{enumerate}
\vspace{2mm}
\noindent\textbf{Q31.} What is the result of "Python"[1:4] ? \hfill [Tag: String basics | Medium]
\begin{enumerate}[label=(\alph*),topsep=2pt,itemsep=1pt,partopsep=0pt,parsep=0pt]
    \item {tho}
    \item {Pyt}
    \item {Python}
    \item {yth}
\end{enumerate}
\vspace{2mm}
\noindent\textbf{Q32.} What is the value of items after executing items = [2, 4, 6]; items.append(8) ? \hfill [Tag: List | Medium]
\begin{enumerate}[label=(\alph*),topsep=2pt,itemsep=1pt,partopsep=0pt,parsep=0pt]
    \item {[2, 4, 8]}
    \item {[2, 4, 6]}
    \item {[8, 2, 4, 6]}
    \item {[2, 4, 6, 8]}
\end{enumerate}
\vspace{2mm}
\noindent\textbf{Q33.} What is the result of int("3.8") ? \hfill [Tag: Type conversions | Medium]
\begin{enumerate}[label=(\alph*),topsep=2pt,itemsep=1pt,partopsep=0pt,parsep=0pt]
    \item {A valid integer conversion}
    \item {4}
    \item {3.8}
    \item {3}
\end{enumerate}
\vspace{2mm}
\noindent\textbf{Q34.} What decimal value does the binary literal 0b1011 represent? \hfill [Tag: Binary numbers | Medium]
\begin{enumerate}[label=(\alph*),topsep=2pt,itemsep=1pt,partopsep=0pt,parsep=0pt]
    \item {13}
    \item {11}
    \item {10}
    \item {9}
\end{enumerate}
\vspace{2mm}
\noindent\textbf{Q35.} Which f-string displays name = "Mira" and score = 92 as Mira scored 92 ? \hfill [Tag: String formatting | Medium]
\begin{enumerate}[label=(\alph*),topsep=2pt,itemsep=1pt,partopsep=0pt,parsep=0pt]
    \item {f"\{name\} scored \{score\}"}
    \item {"\{name\} scored \{score\}"}
    \item {format("name scored score")}
    \item {f"name scored score"}
\end{enumerate}
\vspace{2mm}
\noindent\textbf{Q36.} A file named calculation.py contains print(6 * 7) . Which command runs it from a terminal where Python is available? \hfill [Tag: Using Python shell as a calculator and running simple scripts | Medium]
\begin{enumerate}[label=(\alph*),topsep=2pt,itemsep=1pt,partopsep=0pt,parsep=0pt]
    \item {execute Python calculation.py}
    \item {python calculation.py}
    \item {run calculation.py}
    \item {open calculation.py}
\end{enumerate}
\vspace{2mm}
\noindent\textbf{Q37.} Which function definition follows the usual PEP 8 spacing convention? \hfill [Tag: PEP 8 | Medium]
\begin{enumerate}[label=(\alph*),topsep=2pt,itemsep=1pt,partopsep=0pt,parsep=0pt]
    \item {def add(a , b) :}
    \item {def add (a, b):}
    \item {def add(a, b):}
    \item {def add(a,b):}
\end{enumerate}
\vspace{2mm}
\noindent\textbf{Q38.} Which change most improves the readability of code that calculates a total price? \hfill [Tag: Code readability | Medium]
\begin{enumerate}[label=(\alph*),topsep=2pt,itemsep=1pt,partopsep=0pt,parsep=0pt]
    \item {Remove comments and blank lines}
    \item {Use descriptive names and small functions}
    \item {Place all statements on one line}
    \item {Use shorter names such as x and y}
\end{enumerate}
\vspace{2mm}
\noindent\textbf{Q39.} What does this code print when temperature = 15 ? if temperature $>$ 20: print("warm") else: print("cool") \hfill [Tag: If-else statement | Medium]
\begin{enumerate}[label=(\alph*),topsep=2pt,itemsep=1pt,partopsep=0pt,parsep=0pt]
    \item {cool}
    \item {Nothing}
    \item {warm}
    \item {15}
\end{enumerate}
\vspace{2mm}
\noindent\textbf{Q40.} Which expression evaluates to True ? \hfill [Tag: Equality and relational operators | Medium]
\begin{enumerate}[label=(\alph*),topsep=2pt,itemsep=1pt,partopsep=0pt,parsep=0pt]
    \item {5 == "5"}
    \item {5 != 6}
    \item {5 = 5}
    \item {5 $<$ 3}
\end{enumerate}
\vspace{2mm}
\noindent\textbf{Q41.} A script runs correctly from a terminal using one Python interpreter, but VS Code reports ModuleNotFoundError for a package installed with pip . Which action most directly diagnoses the mismatch? \hfill [Tag: Installing Python and IDEs (Anaconda, Jupyter, VS Code) | Hard]
\begin{enumerate}[label=(\alph*),topsep=2pt,itemsep=1pt,partopsep=0pt,parsep=0pt]
    \item {Compare python -c "import sys; print(sys.executable)" with the selected interpreter}
    \item {Convert the script into a Jupyter notebook}
    \item {Restart VS Code and rerun the script}
    \item {Install the package again without changing interpreters}
\end{enumerate}
\vspace{2mm}
\noindent\textbf{Q42.} What is printed by the following code? \hfill [Tag: Python syntax | Hard]
\begin{enumerate}[label=(\alph*),topsep=2pt,itemsep=1pt,partopsep=0pt,parsep=0pt]
    \item {B}
    \item {A}
    \item {C}
    \item {An indentation error}
\end{enumerate}
\vspace{2mm}
\noindent\textbf{Q43.} After executing the following code, what are the values of a and b ? \hfill [Tag: Variables | Hard]
\begin{enumerate}[label=(\alph*),topsep=2pt,itemsep=1pt,partopsep=0pt,parsep=0pt]
    \item {a == [1, 2, 3] , b == [1, 2]}
    \item {a == [3] , b == [1, 2]}
    \item {a == [1, 2] , b == [1, 2]}
    \item {a == [1, 2, 3] , b == [1, 2, 3]}
\end{enumerate}
\vspace{2mm}
\noindent\textbf{Q44.} Which expression evaluates to a value of type float ? \hfill [Tag: Data types | Hard]
\begin{enumerate}[label=(\alph*),topsep=2pt,itemsep=1pt,partopsep=0pt,parsep=0pt]
    \item {type(7 \% 2)}
    \item {type(7 // 2)}
    \item {type(7 ** 0)}
    \item {type(7 / 2)}
\end{enumerate}
\vspace{2mm}
\noindent\textbf{Q45.} What is the value of -7 // 3 in Python? \hfill [Tag: Operators | Hard]
\begin{enumerate}[label=(\alph*),topsep=2pt,itemsep=1pt,partopsep=0pt,parsep=0pt]
    \item {1}
    \item {-3}
    \item {-2}
    \item {2}
\end{enumerate}
\vspace{2mm}
\noindent\textbf{Q46.} What does the expression 2 ** 3 ** 2 evaluate to? \hfill [Tag: Expressions | Hard]
\begin{enumerate}[label=(\alph*),topsep=2pt,itemsep=1pt,partopsep=0pt,parsep=0pt]
    \item {64}
    \item {256}
    \item {36}
    \item {512}
\end{enumerate}
\vspace{2mm}
\noindent\textbf{Q47.} What is printed when the user enters 08 at the prompt? \hfill [Tag: Input/output | Hard]
\begin{enumerate}[label=(\alph*),topsep=2pt,itemsep=1pt,partopsep=0pt,parsep=0pt]
    \item {81}
    \item {9}
    \item {081}
    \item {A TypeError}
\end{enumerate}
\vspace{2mm}
\noindent\textbf{Q48.} Assume tools.py contains count = 1 and def increment(): global count; count += 1 . What is printed? \hfill [Tag: Module basics | Hard]
\begin{enumerate}[label=(\alph*),topsep=2pt,itemsep=1pt,partopsep=0pt,parsep=0pt]
    \item {1 2}
    \item {2 1}
    \item {1 1}
    \item {2 2}
\end{enumerate}
\vspace{2mm}
\noindent\textbf{Q49.} Which statement correctly describes math.isclose(0.1 + 0.2, 0.3) ? \hfill [Tag: Math module | Hard]
\begin{enumerate}[label=(\alph*),topsep=2pt,itemsep=1pt,partopsep=0pt,parsep=0pt]
    \item {It returns True using tolerance-based comparison}
    \item {It returns False because floating-point values are never equal}
    \item {It converts both operands to integers before comparing}
    \item {It raises an error because decimal literals are approximate}
\end{enumerate}
\vspace{2mm}
\noindent\textbf{Q50.} A UTF-8 encoded file begins with the bytes 72 C3 A9 . Which text does this byte sequence represent? \hfill [Tag: Representing text | Hard]
\begin{enumerate}[label=(\alph*),topsep=2pt,itemsep=1pt,partopsep=0pt,parsep=0pt]
    \item {HÃ(c)}
    \item {He}
    \item {A decoding error}
    \item {eH}
\end{enumerate}
\vspace{2mm}
\noindent\textbf{Q51.} What is the result of the following expression? \hfill [Tag: String basics | Hard]
\begin{enumerate}[label=(\alph*),topsep=2pt,itemsep=1pt,partopsep=0pt,parsep=0pt]
    \item {"Python"}
    \item {"nohtyP"}
    \item {"nohtyp"}
    \item {" nohtyp "}
\end{enumerate}
\vspace{2mm}
\noindent\textbf{Q52.} What is printed? \hfill [Tag: List | Hard]
\begin{enumerate}[label=(\alph*),topsep=2pt,itemsep=1pt,partopsep=0pt,parsep=0pt]
    \item {[2, 3, 9] [2, 3, 9]}
    \item {[1, 2, 3, 4] [1, 2, 3, 4, 9]}
    \item {[1, 2, 3, 4, 9] [2, 3, 9]}
    \item {[1, 2, 3, 4] [2, 3, 9]}
\end{enumerate}
\vspace{2mm}
\noindent\textbf{Q53.} Which statement is correct about the following objects? \hfill [Tag: Common data types summary | Hard]
\begin{enumerate}[label=(\alph*),topsep=2pt,itemsep=1pt,partopsep=0pt,parsep=0pt]
    \item {x is immutable, y is mutable, and z is unordered and unique}
    \item {x is unordered, y is hashable, and z supports integer indexing}
    \item {x , y , and z preserve duplicates and support indexing}
    \item {x and y are immutable, while z preserves insertion duplicates}
\end{enumerate}
\vspace{2mm}
\noindent\textbf{Q54.} What is the result of int(float("3.9")) ? \hfill [Tag: Type conversions | Hard]
\begin{enumerate}[label=(\alph*),topsep=2pt,itemsep=1pt,partopsep=0pt,parsep=0pt]
    \item {3.9}
    \item {4}
    \item {3}
    \item {A ValueError}
\end{enumerate}
\vspace{2mm}
\noindent\textbf{Q55.} What is the decimal value of the Python expression 0b101101 \& 0b011011 ? \hfill [Tag: Binary numbers | Hard]
\begin{enumerate}[label=(\alph*),topsep=2pt,itemsep=1pt,partopsep=0pt,parsep=0pt]
    \item {9}
    \item {45}
    \item {25}
    \item {11}
\end{enumerate}
\vspace{2mm}
\noindent\textbf{Q56.} What is printed by this code? \hfill [Tag: String formatting | Hard]
\begin{enumerate}[label=(\alph*),topsep=2pt,itemsep=1pt,partopsep=0pt,parsep=0pt]
    \item {Ada: 2.3333333333333335}
    \item {name: score}
    \item {Ada: 2.3}
    \item {Ada: 2.33}
\end{enumerate}
\vspace{2mm}
\noindent\textbf{Q57.} A file contains print(\_\_name\_\_) and is executed directly with python file.py . What is normally printed? \hfill [Tag: Using Python shell as a calculator and running simple scripts | Hard]
\begin{enumerate}[label=(\alph*),topsep=2pt,itemsep=1pt,partopsep=0pt,parsep=0pt]
    \item {file.py}
    \item {The absolute file path}
    \item {\_\_main\_\_}
    \item {main}
\end{enumerate}
\vspace{2mm}
\noindent\textbf{Q58.} A function sometimes fails because a caller passes malformed data. Which design best improves maintainability and debugging? \hfill [Tag: Introduction to software development best practices | Hard]
\begin{enumerate}[label=(\alph*),topsep=2pt,itemsep=1pt,partopsep=0pt,parsep=0pt]
    \item {Validate inputs near the function boundary and raise a specific error}
    \item {Catch every exception and return None}
    \item {Duplicate the function to handle each possible input}
    \item {Suppress the failure with a broad try block}
\end{enumerate}
\vspace{2mm}
\noindent\textbf{Q59.} Which revision best follows PEP 8 for a function definition with a long parameter list? \hfill [Tag: PEP 8 | Hard]
\begin{enumerate}[label=(\alph*),topsep=2pt,itemsep=1pt,partopsep=0pt,parsep=0pt]
    \item {def calculate\_total(first, second, third, fourth,):}
    \item {def calculate\_total(first, second, third, fourth):}
    \item {def calculate\_total(first,second,third,fourth):}
    \item {def calculate\_total(first, second, third, fourth):}
\end{enumerate}
\vspace{2mm}
\noindent\textbf{Q60.} Which implementation is most readable and least error-prone for classifying a score? \hfill [Tag: Code readability | Hard]
\begin{enumerate}[label=(\alph*),topsep=2pt,itemsep=1pt,partopsep=0pt,parsep=0pt]
    \item {status = "pass" if score $>$= 60 else "fail"}
    \item {if score $>$= 60: status = "pass"; else: status = "fail"}
    \item {status = "pass" if not score $<$ 60 else "fail"}
    \item {status = ["fail", "pass"][score $>$= 60]}
\end{enumerate}
\vspace{2mm}

\section{Section B: Comprehensive Subjective Problem with Detailed Analytical Rubric}

\begin{subjectiveq}{1}{1}
Design a memory-efficient Python function \texttt{analyze\_numeric\_stream(stream, chunk\_size)} that consumes an iterator of floating-point numbers without loading the entire stream into memory. Calculate the exact mean, sample variance (using Welford's one-pass algorithm), and count of values. Discuss Python object overhead (PyObject structure) versus primitive C types and specify the time and space complexity.
\end{subjectiveq}

\begin{solbox}
\begin{lstlisting}
def analyze_numeric_stream(stream):
    count = 0
    mean = 0.0
    M2 = 0.0
    for x in stream:
        count += 1
        delta = x - mean
        mean += delta / count
        delta2 = x - mean
        M2 += delta * delta2
    if count < 2:
        variance = 0.0
    else:
        variance = M2 / (count - 1)
    return {"count": count, "mean": mean, "variance": variance}
\end{lstlisting}
\textbf{Complexity Analysis:}\\
$\bullet$ \textbf{Time Complexity:} $\mathcal{O}(N)$ where $N$ is the number of stream items, requiring exactly one pass.\\
$\bullet$ \textbf{Space Complexity:} $\mathcal{O}(1)$ auxiliary space since only five scalar variables are retained in the execution frame.\\
\textbf{Memory Architecture Insight:}\\
In CPython, every standard float is a heap-allocated \texttt{PyFloatObject} consuming 24 bytes (8 bytes reference count, 8 bytes type pointer, 8 bytes IEEE 754 value). Consuming data via generators avoids allocating millions of these objects simultaneously.
\end{solbox}

\begin{examinertip}
$\bullet$ \textbf{Rubric Breakdown (10 Marks Total):}\\
1. Algorithmic Formulation \& Theoretical Correctness: \textbf{3 Marks}\\
2. Production Python Syntax \& Idiomatic Standards (PEP 8): \textbf{3 Marks}\\
3. Boundary Case \& Exception Robustness: \textbf{2 Marks}\\
4. Asymptotic Time \& Auxiliary Space Complexity Derivation: \textbf{2 Marks}
\end{examinertip}
